\centerline{\bf \Huge Отбор на подготовительные курсы 7-IT}

\begin{flushright}
        \scriptsize{\textbf{\textit{Именно математика даёт надёжнейшие правила: тому кто им следует — тому не опасен обман чувств. \\Леонард Эйлер}}}
        
\end{flushright}

\problem{\textbf{1}} Начальник записал на доске следующую сумму равную $X$: 

$2025\cdot2024 + 2023\cdot2022 + 2021\cdot2020 + 2019\cdot2018 + 2017\cdot2016 + 2015\cdot2014= X$

Чему равна последняя цифра числа $X$?

\hspace{3mm}

\problem{\textbf{2}} В отборе участвуют только шестиклассники, рассаженные в пять кабинетов. В первом кабинете шестиклассников в два раза меньше, чем во втором, во втором на одного шестиклассника меньше, чем в третьем, в третьем столько же, сколько и в четвертом, а в четвертом на 50\% меньше, чем в пятом. Может ли в пятом кабинете сидеть шестиклассников ровно 36\% от общего числа шестиклассников, участвующих в отборе?

\hspace{3mm}

\problem{\textbf{3}} Чему равно следующее выражение?

$\left(0,025\cdot2000 : \left(0,4 + 1\dfrac{3}{5}\right) + \left(6\dfrac{3}{4}\cdot\dfrac{4}{9} + \dfrac{1}{2} : 0,5\right)\cdot\dfrac{0,6}{0,12}\right):$

$:\dfrac{0,62374-0,17374}{\left(89\cdot98-8650-\dfrac{3\cdot3\cdot3\cdot3\cdot3\cdot3\cdot3}{81}\right)}=?$

\hspace{3mm}

\problem{\textbf{4}} Аркадий и Савр решили устроить лёгкую пробежку вокруг ЭТЛ по часовой стрелке. Каждый раз, когда Аркадий обгоняет Савра, он называет вслух число, равное количеству обгонов им Савра (сразу после первого обгона называет число 1, сразу после второго обгона называет число 2 и так далее). Чему равна сумма всех чисел, названных Аркадием с начала пробежки и до момента, когда он пробежит ровно 2026 кругов вокруг ЭТЛ, если известно, что мальчики начали одновременно с одного места?

\hspace{3mm}

\centerline{\textbf{\textit{Задания 5,6,7 на обратной стороне}}}

\newpage
\hspace{3mm}

\problem{\textbf{5}} Среди 2025 учеников ЭлМШ есть \textit{шиншиллы}, которые всегда говорят правду, и \textit{лососи}, которые всегда лгут (есть хотя бы одна \textit{шиншилла} и хотя бы один \textit{лосось}). Какое наибольшее количество \textit{лососей} может быть, если каждый сказал: "Всего среди 2025 учеников не более 26 \textit{шиншилл}"?

Есть ли логическая ошибка в условии задания? Если есть, объясните её и предложите решение этой проблемы. Если нет, решите задание.

\hspace{3mm}

\problem{\textbf{6}} Однажды Елена решила отправиться из Оргакина (посёлок в Ики-Бурульском районе) в разностороннее путешествие. На своё усмотрение за одно действие она может поехать либо строго на юг, либо строго на север. За первое действие Елена проехала 1км, а за каждое последующее действие в два раза больше предыдущего (2км, 4км, 8км, 16км и так далее). Может ли Елена вернуться в Оргакин за некоторое количество таких действий (Оргакин не объёмный и является точкой на плоскости)?

\hspace{3mm}

\problem{7} Адьян пришел на III ежегодную Математическую Олимпиаду ЭлМШ x ЭТЛ, на которой было предложено 5 задач. Адьян очень торопился и мог потратить на Олимпиаду максимум час. Задачи, придуманные Расулом Владимировичем, Адьян читает за 2 минуты, а решает за 35 минут. Задачи, придуманные Анджуром Аркадьевичем, он читает 2 минуты, а решает за 15 минут. Задачи, придуманные Очиром Владимировичем, он читает 10 минут и решает тоже за 10 минут. Очир Владимирович предложил две задачи, Анджур Аркадьевич одну, а Расул Владимирович две. Какое наибольшее количество баллов Адьян успел набрать, если задачи Очира Владимировича он мог отличить, даже не читая, а отличить задачи Анджура Аркадьевича от задач Расула Владимировича может только прочитав их, при этом, читая, сначала ему попадаются задачи только Расула Владимировича. За задачи Очира Владимировича ученики получали по 4 балла, Анджура Аркадьевича - по 5 баллов, Расула Владимировича - по 8 баллов.

\hspace{3mm}

\centerline{\textbf{\textit{На каждое задание нужно дать развёрнутое решение, а не просто ответ!}}}